\begin{claim}{D.31}{Proved}{Abstract finite groups in the geometric completions of genus at most two}
Let $M_{g,n}^+=\Gamma^+_{g,n}$ allow permutations of the $n$
punctures, and let $M_{g,n}^{\pm}$ also allow orientation reversal.
Every finite nonsolvable subgroup of $\widehat{M_{0,n}^+}$ is
isomorphic to $A_5$. Its permutation action on the punctures has
\[
       n=12a+20b+30c+60m,\qquad
       a,b,c\in\{0,1\},\quad m\geq0,                 \tag{D.31.1}
\]
where the summands are the actual orbit sizes. In particular, it
has at least twelve punctures. Every finite nonsolvable subgroup
of $\widehat{M_{0,n}^{\pm}}$ is isomorphic to $A_5$ or
$A_5\times C_2$. Every finite subgroup of
$\widehat{M_{g,n}^{\pm}}$ for $g=1,2$ is solvable.

These statements determine the abstract groups and, in the positive
sphere case, their puncture-permutation representations. The
surjectivity of the map from geometric $A_5$ conjugacy classes to
completed $A_5$ conjugacy classes, and the completed $A_5$
centralizer comparison of F.7, remain to be proved.
\end{claim}
\begin{proof}
For $n=0,1,2$ the positive sphere mapping class groups are
respectively $1,1,C_2$; including orientation reversal gives
groups of order at most four. Their completions are the same
finite solvable groups. The sphere arguments below therefore
take $n\geq3$.

We give the completed exact sequences first, so that every use of
a subgroup or quotient has its specified completion map.

\emph{A completion and configuration-group preparation.}
For the forgetting sequences used below, the kernel $N$ is finitely
generated and the quotient $Q$ is good. These properties imply
that $\widehat N\to\widehat E$ is injective. Here is the required
argument. Take a characteristic finite-index subgroup $N_0$ of $N$;
these are cofinal because $N$ is finitely generated. In $E/N_0$ the
kernel $B=N/N_0$ is finite. Its centralizer has finite index and fits
into the central extension
\[
 1\longrightarrow Z(B)\longrightarrow C_{E/N_0}(B)
       \longrightarrow Q_0\longrightarrow1,
\]
where $Q_0$ has finite index in $Q$. Goodness passes to $Q_0$.
The class in $H^2(Q_0,Z(B))$ is therefore represented by a continuous
cohomology class of $\widehat Q_0$, hence is inflated from a finite
quotient. Restriction to the kernel of that finite quotient kills
it. The central extension splits over this finite-index subgroup;
a chosen complement meets $B$ trivially and has finite index in
$E/N_0$. Its core is a finite-index normal subgroup still meeting
$B$ trivially. Pulling this core back to $E$ gives a finite quotient
whose restriction to $N$ has kernel exactly $N_0$. This proves the
full induced topology and thus
\[
 1\longrightarrow\widehat N\longrightarrow\widehat E
       \longrightarrow\widehat Q\longrightarrow1.     \tag{D.31.2}
\]
The last equality follows directly from the finite quotients of
$E/N$.

The braid groups on the plane, on finitely punctured surfaces,
and on the closed surfaces of genus one or two used here are
good and have finite cohomological dimension.
We recall how the ingredients apply here. The ordered
configuration-space fibrations of Fadell--Neuwirth
\cite[Theorems~1 and~3]{A-FadellNeuwirth1962} forget one point and
have fibre the surface with finitely many points removed.
The base surface and the fibres are aspherical; hence their
configuration spaces are aspherical by the homotopy exact sequence.
The fibres' groups are finitely generated free groups once a point
has been removed. Starting with the good surface group, the
preceding finite-kernel argument and the Lyndon--Hochschild--Serre
comparison spectral sequences prove goodness inductively.
The required cohomology of each fibre with finite coefficients is
finite. For the plane the initial group is trivial; for the closed
surfaces its goodness is the usual surface-group goodness also
used in \cite{R-BZ}. The unordered group has the ordered one as a
finite-index subgroup, so goodness follows by the finite-index
comparison and Shapiro's lemma. The configuration spaces are
finite-dimensional aspherical manifolds, which give finite
cohomological dimension. Goodness therefore makes each completed
braid group torsion-free: its cohomology with every finite
$p$-primary coefficient module vanishes above that dimension,
whereas a subgroup $C_p$ has infinite $p$-cohomological dimension.
This also follows directly from the prime-subgroup comparison in
\cite[Theorem~A]{R-BZ}.

\emph{Finite subgroups of a completed puncture stabilizer are cyclic.}
For $n\geq4$, capping the boundary of a disc by a once-punctured
disc gives the exact sequence
\[
 1\longrightarrow\langle\Delta^2\rangle\cong\Z
 \longrightarrow B_{n-1}\longrightarrow
 (M_{0,n}^+)_{p}\longrightarrow1.                      \tag{D.31.3}
\]
Here the distinguished puncture is $p$, and the kernel is the
central full twist; this is
\cite[Lemma~2.2]{A-BigelowBudney2001}, whose proof specifies the
capping map. Braid abelianization sends each Artin generator to
$1$ and $\Delta^2$ to $(n-1)(n-2)$. Thus the composite
\[
 \widehat\Z\longrightarrow\widehat B_{n-1}
       \longrightarrow\widehat\Z
\]
is multiplication by this nonzero integer and is injective.
Consequently completion of (D.31.3) is the exact central extension
\[
 1\longrightarrow\widehat\Z\longrightarrow
 \widehat B_{n-1}\longrightarrow
 \widehat{(M_{0,n}^+)_{p}}\longrightarrow1.             \tag{D.31.4}
\]
The last group identifies with the stabilizer of $p$ in
$\widehat{M_{0,n}^+}$: its discrete source has finite index, and
stabilizing $p$ is the inverse image of a subgroup in the finite
puncture-permutation quotient.

Let $K$ be a finite subgroup of that completed stabilizer and let
$E_K$ be its inverse image in $\widehat B_{n-1}$. The subgroup
$A=\widehat\Z$ is central and open in $E_K$, of index $d=|K|$.
Continuous transfer over its finitely many cosets gives
\[
             \operatorname{Ver}:E_K\longrightarrow A,
             \qquad \operatorname{Ver}(a)=d a
             \quad(a\in A).                           \tag{D.31.5}
\]
Its kernel intersects $A$ trivially, so injects into the finite
group $E_K/A$. It is finite, and is trivial because $E_K$ is a
subgroup of the torsion-free completed braid group. Transfer is
therefore a continuous injection into $\widehat\Z$. Its compact
image is a closed procyclic subgroup. Hence $E_K$, and then its
finite quotient $K$, are procyclic and cyclic respectively.
For $n=3$ the same capping map is
$1\to2\Z\to B_2=\Z\to C_2\to1$, and the conclusion is immediate.
For $n=1,2$ the positive point stabilizer is trivial.

It follows in particular that a finite subgroup of a completed
positive sphere mapping class group preserving a specified set
of $1\leq b\leq3$ punctures is solvable. The kernel of its permutation
map to $S_b$ fixes a puncture and is cyclic, while $S_b$ is
solvable. This assertion also applies to every finite-index
decorated subgroup through its actual inclusion into the completed
sphere group.

\emph{The prime-power inputs and elimination of a solvable radical.}
The good-group prime-power theorem
\cite[Theorem~A]{R-BZ} conjugates each finite $p$-subgroup of
$\widehat{M_{0,n}^+}$ to a geometric one and gives its completed
normalizer. The checked prime-power topology argument in F.7
identifies that normalizer with the full completion of its
discrete normalizer. Thus, for a nontrivial normal $p$-subgroup
$A$ of a finite $G\leq\widehat{M_{0,n}^+}$, conjugate $A$ into
$M_{0,n}^+$ and use the exact covering-normalizer sequence
\[
 1\longrightarrow A\longrightarrow N_{M_{0,n}^+}(A)
       \longrightarrow Q_A\longrightarrow1.           \tag{D.31.6}
\]
The group $Q_A$ is a finite-index decorated subgroup of the
mapping class group of the quotient sphere. Its specified branch
set has two points for a cyclic rotation group, and three for a
noncyclic rotation group. This branch set includes the images of
ramification points that were already punctures upstairs; the
ramification indices are retained in the covering decoration.
Completion of (D.31.6) and the normalizer comparison give an
actual injection
\[
                 G/A\lhook\joinrel\longrightarrow\widehat Q_A.
                                                               \tag{D.31.7}
\]
It preserves that set of two or three branch points. The puncture
stabilizer argument makes $G/A$ solvable, and then $G$ is solvable.

For a nontrivial finite solvable normal subgroup $R$ of $G$, the
last nonzero derived term of $R$ is abelian. Its nontrivial
$p$-primary component $A$ is characteristic in $R$, hence normal
in $G$. The preceding argument again makes $G$ solvable.
Consequently a nonsolvable finite $G$ has trivial solvable radical.
Only the prime-power comparison was used in this step.

The same prime-power theorem and the finite rotation-group list
give cyclic odd Sylow subgroups of $G$, and cyclic or dihedral
Sylow $2$-subgroups, including $C_2^2$ in the latter case.
There is also the precise normalizer constraint
\[
 \operatorname{im}\bigl(N_G(C_p)\longrightarrow\Aut(C_p)\bigr)
                      \subset\{1,-1\}.               \tag{D.31.8}
\]
Indeed normalize a geometric rotation. A representative of a
discrete normalizer either fixes its two poles or exchanges them.
The oriented local rotation numbers are opposite, so the induced
automorphism is respectively $1$ or $-1$. The normalizer comparison
preserves this image because $\Aut(C_p)$ is finite. The same
argument with the local orientation sign proves (D.31.8) in the
extended geometric completion as well. Equivalently one can apply
the prime-power normalizer theorem to the good index-two extended
mapping class group; the discrete representatives still induce
only these two exponents.

We record the elementary finite-group reduction needed below.
A finite group all of whose Sylow subgroups are cyclic is solvable.
For the least prime $p$ dividing its order, its cyclic Sylow group
$P$ has $N(P)/C(P)$ of order both prime to $p$ and dividing
$p^{a-1}(p-1)$, where $|P|=p^a$. Minimality of $p$ forces this
quotient to be trivial. Transfer to $P$ has restriction
$x\mapsto x^{[G:P]}$ and is onto. To check this formula without
a fusion assumption, any two elements of the abelian Sylow group
$P$ conjugate in $G$ are conjugate by $N(P)$: the two Sylow
subgroups containing the conjugate element are conjugate inside
its centralizer. Thus they are equal when $P\leq Z(N(P))$.
The coset-cycle formula for transfer now gives the stated exponent.
Its kernel is a normal $p$-complement. Induction on the group order
proves solvability. It follows that every odd-order subgroup here
is solvable, and that a nonsolvable $G$ has a noncyclic dihedral
Sylow $2$-subgroup.

\emph{The verified Gorenstein--Walter application.}
Let $S$ be such a Sylow $2$-subgroup and $z\in Z(S)$ an involution.
The quotient of $C_G(z)$ by $\langle z\rangle$ in (D.31.7)
preserves the two branch points of the rotation quotient.
Its pointwise branch stabilizer is cyclic. Its inverse image
$C_0\leq C_G(z)$ has index at most two and is a central extension
of a finite cyclic group by $\langle z\rangle$, hence is abelian:
a lift of a quotient generator and the central kernel generate it.
The odd-order part of $C_0$ is characteristic in $C_0$, normal in
$C_G(z)$, and cyclic because all odd Sylow subgroups are cyclic.
It is a normal abelian $2$-complement of $C_G(z)$.

We can therefore apply the original, precisely matching theorem
\cite[Theorem~I, p.~553]{A-GorensteinWalter1962}. It states that
$G$ has a normal odd-order subgroup $K$ such that $G/K$ is one
of
\[
       \operatorname{PSL}_2(q),\quad
       \operatorname{PGL}_2(q)\quad(q\text{ odd}),
       \qquad A_7,\qquad\text{a Sylow }2\text{-subgroup}.
                                                               \tag{D.31.9}
\]
The source explicitly includes $C_2^2$ among dihedral groups.
All its hypotheses have just been verified. The odd group $K$
is solvable by the preceding transfer argument, so the triviality
of the solvable radical gives $K=1$.

The group $A_7$ contains the commuting independent cycles
$(1\ 2\ 3)$ and $(4\ 5\ 6)$, contradicting cyclic odd Sylow
subgroups. A $2$-group is solvable. Write $q=p^f$ in either
matrix case. The unipotent subgroup
$u(t)=\left(\begin{smallmatrix}1&t\\0&1\end{smallmatrix}\right)$
for $t\in\F_q$ is elementary abelian of order $p^f$, so $f=1$.
In $\operatorname{PSL}_2(p)$ the diagonal matrices
$\operatorname{diag}(a,a^{-1})$ act on this $C_p$ by
$u(t)\mapsto u(a^2t)$, supplying $(p-1)/2$ distinct automorphisms.
In $\operatorname{PGL}_2(p)$ the matrices
$\operatorname{diag}(a,1)$ supply all $p-1$ automorphisms.
Constraint (D.31.8) leaves $p\leq5$ in the first case and
$p\leq3$ in the second. At $p=3$, the faithful projective-line
action embeds these groups into $S_4$, so they are solvable.
The remaining group is $\operatorname{PSL}_2(5)\cong A_5$.

For this last isomorphism the exact map is conjugation on the
five Sylow $2$-subgroups. The supplied matrix enumeration finds
fifteen involutions, partitions them into the five four-groups,
and verifies that the conjugation image has order sixty and
consists of even permutations. With the recorded ordering of
the five subgroups, the matrices
\[
 \begin{pmatrix}1&1\\0&1\end{pmatrix},\qquad
 \begin{pmatrix}0&-1\\1&0\end{pmatrix}
\]
map respectively to $(1\ 2\ 4\ 3\ 5)$ and
$(2\ 4)(3\ 5)$. Thus the map is injective and its image is exactly
$A_5$, with its sources and targets fixed.

\emph{The puncture representation.}
The pure positive sphere group is the ordered configuration
group of $n-3$ points on the three-punctured sphere: fix the
first three marked points and use their contractible identity
homeomorphism component in the evaluation fibration. It is good
of finite cohomological dimension by the preceding configuration
argument, so its completion is torsion-free. This is exactly
the kernel of the completed map to $S_n$, by the finite-index
completion sequence. Thus every finite subgroup acts faithfully
on the puncture set.

For a finite $A_5$ subgroup, each puncture stabilizer is cyclic by
(D.31.4)--(D.31.5). The cyclic subgroups of $A_5$ have orders
$1,2,3,5$, so its transitive puncture orbits are precisely
$A_5/1$, $A_5/C_2$, $A_5/C_3$, and $A_5/C_5$.
The last three have sizes $30,20,12$. For each $d=2,3,5$,
an element generating $C_d$ fixes
$|N_{A_5}(C_d)/C_d|=2$ points in the corresponding orbit and
none in the other three types. Prime-power conjugacy makes its
puncture permutation that of a rotation of the sphere, which
has at most two fixed punctures. Each exceptional orbit occurs
at most once. This proves (D.31.1), including its actual
permutation representation, and gives the lower bound twelve.

\emph{The extended sphere group.}
Let $G\leq\widehat{M_{0,n}^{\pm}}$ be finite and nonsolvable.
Its positive subgroup $G^+$ has index at most two and is
nonsolvable, so is $A_5$. The centralizer $C_G(A_5)$ intersects
$A_5$ trivially and injects into $G/G^+$, so has order at most
two. With a nontrivial centralizer, its generator is negative
and $G=A_5\times C_2$. With a trivial centralizer, conjugation
embeds $G$ into $\Aut(A_5)\cong S_5$ with image containing the
inner $A_5$, leaving $A_5$ or $S_5$. The automorphism isomorphism
is induced by the action on the five subgroups $A_4\leq A_5$.
An automorphism fixing all five makes $g$ and its image induce
the same permutation of them for every $g\in A_5$; the faithful
natural action of $A_5$ then makes that automorphism the identity.
Conjugation by $S_5$ is faithful as well and supplies one hundred
twenty automorphisms, proving the isomorphism. In $S_5$, the normalizer of
a $5$-cycle acts on it by all four nonzero exponents modulo five.
This violates the extended version of (D.31.8). Therefore only
$A_5$ and $A_5\times C_2$ occur as abstract nonsolvable groups.

\emph{Forgetting punctures in genus two.}
Evaluation on the unordered puncture set, with isotopy extension,
gives the discrete forgetting sequence
\[
 1\longrightarrow B_n(S_2)\longrightarrow M_{2,n}^+
        \longrightarrow M_{2,0}^+\longrightarrow1.      \tag{D.31.10}
\]
The identity homeomorphism component of $S_2$ is contractible,
as used in Part C, so there is no preceding evaluation-loop
kernel. The quotient is good, the braid kernel is finitely
generated, and (D.31.2) completes this sequence exactly.
The completed kernel is torsion-free by the configuration
argument. Every finite subgroup of $\widehat{M_{2,n}^+}$
therefore injects under forgetting into $\widehat{M_{2,0}^+}$.
The hyperelliptic quotient is the exact sequence
\[
 1\longrightarrow C_2\longrightarrow M_{2,0}^+
        \longrightarrow M_{0,6}^+\longrightarrow1;
                                                               \tag{D.31.11}
\]
see \cite[Propositions~3.2--3.3]{A-BigelowBudney2001}.
The central involution is detected on homology modulo three,
so it injects into the completion. Finite-kernel completion
then makes (D.31.11) exact with hats. Every finite subgroup
has solvable image in $\widehat{M_{0,6}^+}$ by (D.31.1),
and kernel contained in $C_2$. It is solvable. The same holds
in the extended group because the positive subgroup has
index at most two. The case $n=0$ uses (D.31.11) directly.

\emph{Forgetting punctures in genus one.}
For $n\geq1$, evaluation gives
\[
 1\longrightarrow N_n\longrightarrow M_{1,n}^+
       \longrightarrow\operatorname{SL}_2(\Z)
       \longrightarrow1,\qquad
 N_n=B_n(T^2)/\Z^2.                                    \tag{D.31.12}
\]
Here $\Z^2$ is the image of the translation loops in the
identity homeomorphism component, identified by evaluation
as in C.9. It is central in the braid group: applying a loop
of translations to every configuration of a braid gives the
square homotopy interchanging the two loop products. Its injection
can be checked on the ordered configuration space by evaluating
the first strand. To verify its full topology in the unordered
braid group, sum the homology classes of all strand paths:
\[
 B_n(T^2)\longrightarrow H_1(T^2,\Z)=\Z^2.
\]
Endpoints are permuted, so the sum is a cycle and gives a
homomorphism. On diagonal translation loops it is multiplication
by $n$. The resulting completed injection proves the exact
central extension
\[
 1\longrightarrow\widehat\Z^2\longrightarrow
 \widehat{B_n(T^2)}\longrightarrow\widehat N_n
       \longrightarrow1.                              \tag{D.31.13}
\]
The completed braid group is torsion-free by the configuration
argument. For completeness the ordered torus configuration
has the explicit product coordinates
\[
 (x_1,\ldots,x_n)\longmapsto
 (x_1,x_2-x_1,\ldots,x_n-x_1)
 \in T^2\times\operatorname{Conf}_{n-1}(T^2\setminus\{0\}),
\]
so its goodness also follows directly from that of $\Z^2$
and the iterated punctured-surface configuration groups.

For a finite $K\leq\widehat N_n$, pull it back in (D.31.13).
Transfer to the central open $\widehat\Z^2$ has restriction
$a\mapsto |K|a$. Its kernel again injects into $K$ and is
finite, so torsion-freeness makes transfer injective.
The pullback, and hence $K$, is abelian. Since
$\operatorname{SL}_2(\Z)$ is good and $N_n$ is finitely
generated, (D.31.2) also completes (D.31.12) exactly.
The image of every finite subgroup lies in a finite subgroup
of $\widehat{\operatorname{SL}_2(\Z)}$, which is cyclic
of order $1,2,3,4$, or $6$ by D.8's virtually free comparison.
The finite kernel is abelian by (D.31.13), proving solvability.
For orientation reversal the positive subgroup has index at
most two, which preserves this conclusion. For $n=0$ use
D.8 directly with $\operatorname{GL}_2(\Z)$.

The executed program
\path{code/part_a/sphere_completion_restriction.py}
checks the projective-matrix automizer calculations at
$p=3,5,7,11,13$, constructs the displayed
$\operatorname{PSL}_2(5)\to A_5$ isomorphism, compares
the $C_5$ normalizer actions in $A_5$ and $S_5$, and
checks the exceptional coset fixed-point counts and the
$C_3^2$ subgroup of $A_7$. The general reductions are
the completed exact sequences and transfer arguments above.
\end{proof}
